% Independently derived directly from the defining complete elliptic integrals.
% This file is an inclusion, and requires only standard mathematical macros.
\subsection{Exact derivative and uniform bounds for the elliptic profile}

For $0<\rho<1$, retain the original modulus
$k=\sqrt{1-\rho^2}$ and define
\[
 D_\rho(\theta)=\cos^2\theta+\rho^2\sin^2\theta,
 \qquad
 K=K(k)=\int_0^{\pi/2}\frac{d\theta}{\sqrt{D_\rho(\theta)}},
 \qquad
 E=E(k)=\int_0^{\pi/2}\sqrt{D_\rho(\theta)}\,d\theta.
\]
Put
\[
 A=\int_0^{\pi/2}\frac{\sin^2\theta}{\sqrt{D_\rho(\theta)}}\,d\theta,
 \qquad
 B=\int_0^{\pi/2}\frac{\cos^2\theta}{\sqrt{D_\rho(\theta)}}\,d\theta,
 \qquad \mathfrak q(\rho)=\frac{\rho K}{E}.
\]
All these functions are $C^\infty$ on $(0,1)$: on each compact
subinterval, $D_\rho$ has a positive lower bound, and every derivative
of each integrand has a uniform integrable bound. Directly,
\[
 A+B=K,\qquad K-E=(1-\rho^2)A,
 \qquad E-\rho^2K=(1-\rho^2)B,
 \qquad E_\rho=\rho A.
\]
The function $\sin\theta\cos\theta/\sqrt{D_\rho(\theta)}$
vanishes at both integration endpoints and has derivative
\[
 \frac{\cos^4\theta-\rho^2\sin^4\theta}
 {D_\rho(\theta)^{3/2}}.
\]
Its integral is zero. Since
\[
 \rho^2\sin^2\theta-\cos^2\theta D_\rho(\theta)
 =\rho^2\sin^4\theta-\cos^4\theta,
\]
this proves
\[
 \int_0^{\pi/2}\frac{\sin^2\theta}{D_\rho(\theta)^{3/2}}\,d\theta
 =\frac B{\rho^2},\qquad K_\rho=-\frac B\rho.
\]
Consequently the derivative relevant to the implicit equation is exactly
\[
 \boxed{
 \mathfrak q'(\rho)=\frac{K-B}{E}-\frac{\rho^2KA}{E^2}
 =\frac{(1-\rho^2)AB}{E^2}
 =\frac{(K-E)(E-\rho^2K)}{(1-\rho^2)E^2}>0.}
\]
In particular this derivative never vanishes at an interior parameter.
As $\rho\uparrow1$, dominated convergence gives $K,E\to\pi/2$
and $\mathfrak q(\rho)\to1$. To check the other endpoint without an elliptic
asymptotic formula, Cauchy--Schwarz and $u=\tan\theta$ give
\[
 K^2\le\frac\pi2\int_0^{\pi/2}\frac{d\theta}{D_\rho(\theta)}
 =\frac\pi2\int_0^\infty\frac{du}{1+\rho^2u^2}
 =\frac{\pi^2}{4\rho}.
\]
Also $E\ge\int_0^{\pi/2}\cos\theta\,d\theta=1$.
Thus $0<\mathfrak q(\rho)\le(\pi/2)\sqrt\rho\to0$ as $\rho\downarrow0$.
Strict monotonicity and these two limits show that $\mathfrak q$ is a
$C^\infty$ bijection from $(0,1)$ to $(0,1)$ with $C^\infty$ inverse.
Therefore the original equation
\[
 \mathfrak q(\rho(t))=\frac1{1+t},\qquad t>0,
\]
defines a unique $C^\infty$ function $\rho(t)$, strictly decreasing,
and $f(t)=\log((1+\rho(t))/(1-\rho(t)))$ is $C^\infty$.
Differentiating the exact equation gives
\[
 \rho'(t)=-\frac{\mathfrak q(\rho(t))^2}{\mathfrak q'(\rho(t))},\qquad
 f'(t)=-\frac{2\mathfrak q(\rho(t))^2}{(1-\rho(t)^2)\mathfrak q'(\rho(t))}.
\]
Using $t=(1-\mathfrak q)/\mathfrak q$, this becomes
\[
 \boxed{
 t f'(t)=-\frac{2\rho K(E-\rho K)}{(K-E)(E-\rho^2K)}<0.}
\]
Every $K,E,\rho$ in this expression has its original value at $t$.

Here is a uniform bound requiring no numerical elliptic evaluation.
Set $e_\rho=E/K$. The already proved inequality $\mathfrak q<1$ gives $e_\rho>\rho$.
For any positive $x,y$, $(x-y)(x^{-1/2}-y^{-1/2})\le0$.
Apply this with $x=D_\rho(\theta)$ and $y=D_\rho(\phi)$ and
integrate over $[0,\pi/2]^2$. Expanding the four terms yields
\[
 2\frac\pi2 E-2K\int_0^{\pi/2}D_\rho(\theta)\,d\theta\le0.
\]
Since the last integral is $\pi(1+\rho^2)/4$, this proves
\[
 \rho<e_\rho\le\frac{1+\rho^2}{2}<1.
\]
The polynomial $p(x)=(1-x)(x-\rho^2)$ has derivative
$1+\rho^2-2x\ge0$ on $[\rho,(1+\rho^2)/2]$.
It follows that
\[
 (1-e_\rho)(e_\rho-\rho^2)\ge p(\rho)=\rho(1-\rho)^2,
 \qquad 2\rho(e_\rho-\rho)\le\rho(1-\rho)^2.
\]
Substitution in the exact derivative proves
\[
 \boxed{0<-t f'(t)
 =\frac{2\rho(e_\rho-\rho)}{(1-e_\rho)(e_\rho-\rho^2)}\le1\qquad(t>0).}
\]
In particular, for any $0<t_1<t_2$, the fundamental theorem of calculus
gives the finite-difference estimate
\[
 0<f(t_1)-f(t_2)\le\log(t_2/t_1).
\]

For completeness, all constants in the small-$t$ assertion can also
be bounded directly from these integrals. Since $D_\rho\le1$,
$A,B\ge\pi/4$ and $E\le\pi/2$, whence $AB/E^2\ge1/4$.
On the other hand, $AB\le(A+B)^2/4=K^2/4$ and $E>\rho K$,
whence $AB/E^2\le1/(4\rho^2)$. Thus
\[
 \frac{1-\rho^2}{4}\le \mathfrak q'(\rho)
 \le\frac{1-\rho^2}{4\rho^2}.
\]
Write $d_\rho=1-\rho$ and $b_\rho=1-\mathfrak q(\rho)$.
Integration from $\rho$ to $1$ gives the exact bounds
\[
 \frac{d_\rho^2(3-d_\rho)}{12}\le b_\rho
 \le\frac{d_\rho^2}{4\rho}.
\]
Also $K\ge E$, so $\mathfrak q\ge\rho$ and
\[
 \frac{d_\rho^2(3-d_\rho)}{12}\le t=\frac{b_\rho}{\mathfrak q}
 \le\frac{d_\rho^2}{4\rho^2}.
\]
If $\rho\le1/2$, the lower bound for $b_\rho$ at $\rho=1/2$ and
monotonicity imply $t\ge b_\rho\ge5/96>1/20$.
Hence $0<t\le1/20$ implies $0<d_\rho<1/2$ and
\[
 d_\rho\le\sqrt{24/5}\,\sqrt t.
\]
Set $W=t(1+\rho)^2/d_\rho^2$. The preceding bounds give
\[
 W\ge\frac{(2-d_\rho)^2(3-d_\rho)}{12}
 =1-\frac43d_\rho+\frac7{12}d_\rho^2-\frac1{12}d_\rho^3
 \ge1-\frac43d_\rho,
 \qquad
 W\le\left(\frac{1+\rho}{2\rho}\right)^2.
\]
Using $\log(1+x)\le x$ for $x\ge0$ and
$\log(1-x)\ge-x/(1-x)$ for $0\le x<1$, with
$d_\rho\le1/2$, yields
\[
 -2d_\rho\le\frac12\log W
 =f(t)+\frac12\log t\le d_\rho.
\]
Thus, with its original constant term retained,
\[
 \boxed{
 \left|f(t)+\frac12\log t\right|
 \le2\sqrt{24/5}\,\sqrt t<\frac{22}{5}\sqrt t
 \qquad(0<t\le1/20).}
\]

To obtain the derivative remainder independently of differentiating
an $O$-term, set $T=-tf'(t)=2\mathfrak q b_\rho/((1-\rho^2)\mathfrak q')$.
The same bounds give
\[
 T\ge\frac{2\rho^3(2+\rho)}{3(1+\rho)^2}
 \ge\frac{\rho^3}{2},
 \qquad T\le\frac{2}{\rho(1+\rho)^2}.
\]
The second lower inequality follows from
$4(2+\rho)-3(1+\rho)^2=(1-\rho)(3\rho+5)\ge0$.
For $1/2\le\rho<1$, the lower estimate implies
$T-1/2\ge-3d_\rho/2$, since $(1-d_\rho)^3\ge1-3d_\rho$.
For the upper estimate,
\[
 \frac{2}{\rho(1+\rho)^2}-\frac12
 =\frac{d_\rho(\rho^2+3\rho+4)}{2\rho(1+\rho)^2}
 \le\frac{32}{9}d_\rho<4d_\rho.
\]
Consequently
\[
 \boxed{
 \left|t f'(t)+\frac12\right|
 \le4\sqrt{24/5}\,\sqrt t<9\sqrt t
 \qquad(0<t\le1/20),}
\]
equivalently $f'(t)=-1/(2t)+O(t^{-1/2})$ with the displayed
explicit constant. These arguments prove both
$f(t)=-\tfrac12\log t+O(\sqrt t)$ and its derivative assertion
directly, without differentiating an unproved remainder estimate.

Finally, when $0<t\le4$, one has $\mathfrak q\ge1/5$; the previously proved
$\mathfrak q\le(\pi/2)\sqrt\rho$ gives
$\rho\ge4/(25\pi^2)$. The lower bound $T\ge\rho^3/2$
therefore supplies, if a strictly positive uniform lower constant is needed,
\[
 \boxed{\frac{32}{15625\pi^6}\le -t f'(t)\le1
 \qquad(0<t\le4).}
\]


